\section{Scope, Chronology, and Evidence}
\label{sec:aquinas-apparatus}

\subsection{Identity and bounds}

\textbf{Name and life.} This study uses Saint Thomas Aquinas, Latin Thomas de
Aquino: Dominican friar, priest, theologian, philosopher, biblical commentator,
saint, and Doctor of the Church.  He was born about 1225, within a responsible
1224--1226 reconstruction, and died at Fossanova on 7 March 1274.  No exact
birthday is claimed.  ``Angelic Doctor'' and ``Common Doctor'' are received
titles, not contemporary civil identifiers.

\textbf{Places and institutions.} The principal route is Roccasecca and Aquino;
Monte Cassino; Naples; Paris; Cologne; the papal-court cities of central Italy,
especially Orvieto; Rome; Naples again; Maenza; and Fossanova.  Modern place
names orient without deciding disputed medieval boundaries.  Governing
institutions include the Aquino household, Benedictine Monte Cassino, the Order
of Preachers, the universities and studia of Naples, Paris, Cologne, and Rome,
the papal court, and Cistercian Fossanova.

\textbf{Language and ecclesial frame.} Thomas wrote principally in Latin and
read Greek, Arabic, Hebrew, and Jewish authorities largely through Latin
translations, excerpt collections, and scholarly mediation.  The work is a
historical and hagiographic biography within Latin Catholic history.  It does
not make the Roman Church the only community capable of reading Thomas or
treat thirteenth-century Latin formulations as the only Christian vocabulary.

\textbf{Currentness.} Official, institutional, digital-corpus, liturgical,
relic-custody, and bibliographic claims retained here were checked through 20
July 2026.  The current Roman memorial is 28 January.  The present Toulouse
custodial claim and the 2025 co-patronage formulation are mutable and must be
rechecked in a future revision.  Calendar claims are limited to the General
Roman Calendar; Dominican, local, other-rite, and historical calendar editions
were not comprehensively audited.

\subsection{Reader, question, and thesis boundary}

The intended reader is a serious general reader, cleric, student, or researcher
who wants a sustained life before a summary of doctrines.  The governing
question is how Thomas's Dominican vocation, schools, collaborators,
controversies, prayer, and unfinished work produced a durable account of
wisdom, and how that achievement can be received without flattening evidence
or historical limits.

The historical thesis is that Thomas made distinction serve communion: divine
and created causality, faith and reason, body and soul, nature and grace,
contemplation and teaching remain really distinguishable precisely so they can
be rightly ordered.  This architecture grew from a corporate religious and
academic life, not solitary abstraction.  It also contains positions shaped by
inadequate sources and medieval coercive, anti-Jewish, gendered, and servile
orders.  Later sanctity and Doctor judgments are stated according to their
ecclesial objects; this study does not issue those judgments, revoke them, or
extend them to every proposition and anecdote.

Included are the controlled life sequence; principal writing genres and
development; Dominican, university, papal, court, and collaborative settings;
representative theological and philosophical structures; the mendicant and
Aristotelian controversies; difficult moral and social conclusions;
cessation, death, burial, canonization, relic conflict, cult, iconography,
patronage, and modern Catholic reception.  Excluded are a full intellectual
commentary, complete work and manuscript catalog, exhaustive Thomism history,
canonization dossier edition, forensic relic judgment, and independent
theological or ecclesiastical verdict.

\subsection{Evidence key}

\begin{evidencelegend}
\evidencelegendrow{A}{Subject-authored or authorial evidence: Thomas's works,
autograph material, letters, and author-supervised texts, controlled by genre,
chronology, textual status, and transmission.  A work is not private memoir.}
\evidencelegendrow{N}{Contemporary or near-contemporary external evidence:
charters, university and papal acts, early catalogs, participant or companion
testimony.  A canonization deposition preserves different proximity levels
inside one official record.}
\evidencelegendrow{E}{Early received report: especially William of Tocco,
Bernard Gui, Peter Calo, Ptolemy of Lucca, and hearsay in the inquiries.  These
may preserve lost memory while arranging it for sanctity and reception.}
\evidencelegendrow{L}{Later hagiographic, liturgical, artistic, scholastic,
relic, quotation, and patronage tradition; evidence first for the age that
transmits it.}
\evidencelegendrow{M}{Official ecclesial reception, limited to the act's
object: canonization, Doctor declaration, feast, patronage, conciliar teaching,
or papal recommendation.}
\evidencelegendrow{K}{Named modern critical reconstruction, including work
chronology, source dependence, condemnations, social history, and relic
custody, with its premises and alternatives retained.}
\evidencelegendrow{S}{Source-grounded project synthesis; never attributed to
Thomas or to an ecclesial act.}
\end{evidencelegend}

\subsection{Detailed chronology}

\begin{biotimeline}
\biotimelinerow{1224--1226}{Roccasecca / Aquino region}{Birth to Landulf and
Theodora in an aristocratic household.}{K range; no birthday}
\biotimelinerow{c.~1230s}{Monte Cassino}{Childhood education in relation to the
Benedictine abbey.}{E with documentary family context}
\biotimelinerow{c.~1239--1244}{Naples}{Study at Frederick II's studium;
encounter with Aristotelian curriculum and the Order of Preachers.}{E/K;
curriculum reconstructed}
\biotimelinerow{c.~1244}{Naples / road north}{Entry into the Dominicans;
family members intercept him during removal from Naples.}{E; sequence strong}
\biotimelinerow{c.~1244--1245}{Family strongholds}{Detention and continued
refusal to abandon the order; eventual return to Dominican direction.}{E/K;
duration and scenes variable}
\biotimelinerow{1245--1248}{Paris}{Dominican theological and philosophical
formation under Albert's influence.}{Institutional chronology / E}
\biotimelinerow{1248--1252}{Cologne}{Studies with Albert at the Dominican
studium.}{E/K; classroom anecdotes later}
\biotimelinerow{1248--1252}{Cologne}{Probable priestly ordination during the
Cologne years.}{K; exact date and prelate unresolved}
\biotimelinerow{1252--1256}{Paris}{Bachelor in Scripture and the
\emph{Sentences}; early works including \emph{On Being and Essence}.}{Work and
university chronology}
\biotimelinerow{1256--1259}{Paris}{Licence and first regency as master;
secular--mendicant controversy and disputed questions.}{Documents / A / K}
\biotimelinerow{1259}{Valenciennes / Italy}{Associated with Dominican study
organization; begins a multi-city Italian assignment.}{General-chapter record;
individual contribution bounded}
\biotimelinerow{c.~1259--1265}{Italy}{Composition of much of the \emph{Summa
contra gentiles}; purpose and addressee debated.}{A manuscript/work chronology
/ K}
\biotimelinerow{1261--1265}{Orvieto / papal court}{\emph{Catena aurea}, Greek
dossier work, preaching, and early Corpus Christi attribution under Urban IV.}{A/E;
attribution qualified}
\biotimelinerow{c.~1265--1268}{Rome}{Dominican teaching and beginning of the
\emph{Summa theologiae}.}{A/K; exact institutional title debated}
\biotimelinerow{1268/1269--1272}{Paris}{Second regency; mendicant polemic,
Aristotelian commentaries, intellect and creation disputes, quodlibets, and
continued \emph{Summa}.}{A/documents/K}
\biotimelinerow{10 Dec.~1270}{Paris}{Tempier condemns thirteen propositions in
the intellectual controversy; Thomas is not named.}{M/documentary act}
\biotimelinerow{1272--1273}{Naples}{Teaches theology, preaches, and continues
Part III; a royal order funds his lectures.}{N/A/E/K}
\biotimelinerow{about 6 Dec.~1273}{Naples}{After Mass, Thomas undergoes a
reported change and ceases ordinary composition; the \emph{Summa} remains
unfinished at III, q.~90.}{Early indirect report / corpus fact}
\biotimelinerow{early 1274}{Road toward Lyon}{Begins journey after summons to
the Second Council of Lyon; illness worsens.}{E; accident details qualified}
\biotimelinerow{7 Mar.~1274}{Fossanova}{Death and burial among the
Cistercians.}{Early chronological convergence}
\biotimelinerow{7 Mar.~1277}{Paris}{Tempier condemns 219 propositions; Thomas
is unnamed, but some articles implicate positions associated with him.}{M/K}
\biotimelinerow{18 Mar.~1277}{Oxford}{Kilwardby prohibits theses including
positions more directly associated with Thomas.}{M/K; scope local}
\biotimelinerow{1318--1323}{Italy / Avignon}{Tocco promotes the cause and
revises his \emph{Ystoria}; inquiries at Naples (1319) and Fossanova (1321)
gather testimony.}{Official dossier / E}
\biotimelinerow{18 Jul.~1323}{Avignon}{John XXII canonizes Thomas.}{M}
\biotimelinerow{1325}{Paris}{The 1277 act is revoked only insofar as it touches
or is asserted to touch the canonized Thomas.}{M; no blanket doctrinal approval}
\biotimelinerow{1368--1369}{Fossanova / Toulouse}{Principal relic translation
to Dominican custody at Toulouse after prolonged conflict.}{M/E/K; body
history divided}
\biotimelinerow{11 Apr.~1567}{Rome}{Pius V declares Thomas a Doctor of the
Church in \emph{Mirabilis Deus}.}{M}
\biotimelinerow{1879--1880}{Rome}{\emph{Aeterni Patris} promotes Thomistic
renewal; \emph{Cum hoc sit} names Thomas patron of Catholic schools.}{M}
\biotimelinerow{1974}{Rome}{Paul VI's \emph{Lumen Ecclesiae} marks the seventh
centenary and states both Thomistic preference and legitimate plurality.}{M}
\biotimelinerow{1 Nov.~2025}{Rome}{Leo XIV describes Aquinas and Newman as
co-patrons of the Church's educational mission.}{M; mutable wording}
\end{biotimeline}

\subsection{Method, source limits, and negative results}

The narrative coordinates Thomas's authenticated corpus, contemporary
documents, the \emph{Fontes vitae}, official reception, and named modern
reconstruction.  Canonization depositions are read witness by witness: a
participant report is not merged with hearsay merely because both occur in the
same proceeding.  Tocco is both a crucial collector and the cause's promoter.
Gui and Calo are not counted as independent corroboration where they depend on
the same tradition.  Work dates follow critical ranges and never create a
private psychological diary from a proposition.

No controlled source established an exact birth date; complete school
transcript; exact detention duration; independent contemporary witness for the
``dumb ox,'' temptation, angelic girdle, or crucifix speech; Reginald's own
deposition for the straw report; a clinical diagnosis
or poisoning evidence; one untouched relic chain; contemporary portrait; or
universal present patronage beyond named official acts.  Absence within this
bounded search is not proof that no local record exists.

The study inspected representative primary loci for doctrine and the difficult
legacy.  It is not a complete collation of the Leonine edition, every
reportatio, all canonization witnesses, the Paris and Oxford dossiers, local
archives, relic inventories, or Thomistic schools.  Greek, Arabic, Hebrew, and
Jewish sources were not independently collated against every Latin form Thomas
used.

\subsection{Rights and review state}

Project-created prose and organization are intended for CC BY 4.0.  Medieval
Latin wording is public domain, but modern critical texts, editorial apparatus,
translations, scholarship, web presentation, manuscript images, Vatican pages,
and institutional catalog material retain their own status.  Corpus
Thomisticum's electronic presentation is expressly rights-reserved and is
linked and paraphrased, not reproduced.  The \emph{Fontes vitae} scans were
used as page-image and bibliographic controls; OCR is only a finding aid.  No
third-party image or extended protected translation is reproduced.

Internal source, chronology, tradition, metadata, production, and page review
are recorded for this working study.  Outstanding are independent medievalist,
Dominican, Thomist, philosophical, theological, Jewish-Christian,
gender-history, slavery-history, liturgical, canonization-process, relic,
rights, and ecclesiastical review; a complete work and manuscript collation;
and a work-specific rights review for distribution.  The publication has no
imprimatur, nihil obstat, or ecclesiastical approval and remains a study aid.
